\documentclass[11pt,a4paper]{article}
\usepackage[margin=19mm]{geometry}
\usepackage[T1]{fontenc}
\usepackage{lmodern,graphicx,booktabs,amsmath,xcolor,hyperref,fancyhdr}
\hypersetup{colorlinks=true,urlcolor=blue,linkcolor=blue,pdftitle={Deep hedging: training and evaluation curves}}
\pagestyle{fancy}\fancyhf{}\lhead{Deep hedging experiments}\rhead{18 October 2026}\cfoot{\thepage}
\setlength{\headheight}{14pt}
\setlength{\parindent}{0pt}\setlength{\parskip}{7pt}
\newcommand{\plot}[2]{\begin{center}\includegraphics[width=\linewidth,height=#2,keepaspectratio]{figures/#1.pdf}\end{center}}
\begin{document}
\begin{center}
{\LARGE\bfseries Deep hedging}\\[5pt]
{\Large Training losses and evaluation curves}\\[8pt]
Results snapshot: 18 October 2026
\end{center}

\section*{Experiment and main observations}
This report reconstructs learning curves from the Hetzner archive:\\
\texttt{/home/jump/deep-hedging/results}. It covers the seven scientific
\texttt{-lr001} runs, excluding smoke and profiling runs. Each method completed
2,700 seconds (45 minutes) of active training with seed 20261017.

The selected fixed base-distribution baseline has the lowest final base MSE,
easy MSE and balanced scaled loss. SOAR mixed has the lowest hard MSE, improving
on baseline by approximately 0.30\%, while worsening its base and easy losses.
These are single-seed observations, not statistically established rankings.

\subsection*{What is being measured?}
The per-episode loss is squared normalized hedging shortfall:
\[
 Z=\frac{\text{payoff}-\text{initial premium}-\text{trading gains}
                 +\text{transaction and liquidation costs}}{S_{\mathrm{inception}}},
 \quad L_d=\frac{1}{N_d}\sum_{i=1}^{N_d} Z_i^2.
\]
The CSV field \texttt{raw\_mse} means $L_d$ before reference-scale normalization;
it is already spot-normalized and is not a dollar-squared loss. Lower is better.
The development checkpoint-selection objective is
\[
 B=\tfrac12\left(\frac{L_{\mathrm{easy}}}{s_{\mathrm{easy}}}
                  +\frac{L_{\mathrm{hard}}}{s_{\mathrm{hard}}}\right),
 \quad s_{\mathrm{easy}}=0.00105552154,\quad
 s_{\mathrm{hard}}=0.00281462749.
\]
Base is reported separately, with reference scale $s_{\mathrm{base}}=0.00150060048$.

\begin{center}\begin{tabular}{lll}\toprule
Data stream & Episodes per domain & Purpose\\\midrule
Training & 256 per learner step & Fresh samples from each method's generator\\
Monitor & 512 & Fixed evaluation data during training\\
Selection & 1,024 & Fixed development data for checkpoint choice\\
Final test & 8,192 & Evaluation after selection was frozen\\\bottomrule
\end{tabular}\end{center}

Training logs record the mean of four pre-update minibatch losses in a sampled
block, roughly every 30 active seconds. They are not averages over the intervening
30 seconds. The report uses 91 records per method, including the screening and
final terminal records. Active time excludes compilation, evaluation and
checkpoint/history I/O; it is not total paid wall time.

The following pages separate sampled training loss, monitor curves, development
selection and retrospective final-test curves. Lines join measured checkpoints;
no evaluation smoothing or uncertainty bands are applied.

\clearpage
\section*{1. Loss during training}
\plot{training}{0.76\textheight}
{\small Faint lines show all logged block losses; bold lines show a centered
five-record rolling median (fewer records at the edges). All panels share a
logarithmic vertical scale. The smoother is descriptive only.}

The hard reference samples a fixed hard distribution; adaptive methods change
their training distribution. Thus a higher training loss can reflect harder
sampled episodes. These training curves alone do not rank generalization.

\clearpage
\section*{2. Monitor evaluation during training}
\plot{monitor}{0.75\textheight}
{\small Fixed monitor set: 512 episodes per domain. Left: the full trajectory,
with a logarithmic MSE axis. Right: a linear-axis zoom after five active minutes.
The plotted quantity is $10^4L_d$; each row has its own vertical range.}

The full view exposes the initial improvement, while the zoom resolves later
checkpoint fluctuations and differences between methods. Domain-specific raw
loss levels should be compared across methods within the same domain. Monitor
results are distinct from the development selection objective.

\clearpage
\section*{3. Development evaluation and checkpoint selection}
\plot{selection}{0.66\textheight}
{\small Fixed selection set: 1,024 episodes per domain. Curves show recorded
checkpoints after five minutes; the untrained checkpoint is omitted from this
zoom. Open rings in the balanced-loss panel identify the frozen choices. Base
loss does not enter the balanced easy/hard selection objective.}

Selection was frozen on 18 October 2026 at 01:54:09 UTC, before final-test
evaluation. Six methods selected their terminal 45-minute checkpoint; SOAR mixed
selected its 30-minute checkpoint (1,800.06 active seconds), although its lineage
continued to 45 minutes. The final table therefore compares development-selected
checkpoints, not necessarily the last checkpoint of every method.

\clearpage
\section*{4. Retrospective final-test curves}
\plot{final_test}{0.75\textheight}
{\small Fixed final-test set: 8,192 episodes per domain. These evaluations were
performed after selection was frozen. The horizontal axis denotes the training
age of the saved checkpoint, not the time when its test evaluation ran. Left:
full trajectory, logarithmic MSE. Right: linear zoom after five minutes.}

These curves describe how saved policies perform on the final datasets. They
were not the online feedback stream and must not be used to retroactively change
the reported checkpoint selection. Repeated checkpoints share the same dataset;
they are not independent replications.

\clearpage
\section*{5. Selected-checkpoint results}
\begin{center}\small
\begin{tabular}{lrrrrr}\toprule
Method & Time (min) & Base & Easy & Hard & Balanced\\
 & & \multicolumn{3}{c}{MSE $\times 10^4$} & scaled loss\\\midrule
Baseline & 45.0 & 0.6657 & 0.4248 & 1.9968 & 0.055596 \\
Hard reference & 45.0 & 0.7213 & 0.4685 & 1.9948 & 0.057629 \\
DPG hard & 45.0 & 0.7262 & 0.4629 & 2.0320 & 0.058026 \\
SOAR hard & 45.0 & 0.6811 & 0.4315 & 1.9949 & 0.055879 \\
Selfplay & 45.0 & 0.7172 & 0.4677 & 2.0321 & 0.058255 \\
DPG mixed & 45.0 & 0.7161 & 0.4561 & 2.0071 & 0.057260 \\
SOAR mixed & 30.0 & 0.6848 & 0.4372 & 1.9908 & 0.056075 \\
\bottomrule
\end{tabular}\end{center}

The baseline has the best base, easy and balanced test losses. SOAR mixed reduces
hard MSE by $5.96\times10^{-7}$ relative to baseline, but increases base MSE by
$1.91\times10^{-6}$ and easy MSE by $1.24\times10^{-6}$. Hard reference and SOAR
hard have still smaller measured hard-domain advantages. Neither DPG variant nor
selfplay improves any of the three selected-test MSEs over baseline.

\plot{diagnostics}{0.38\textheight}
{\small Left: learner updates versus active time. Right: logged generator KL
relative to the base distribution, on a symmetric-log scale with a linear region
near zero. These diagnostics contextualize loss trajectories; they do not isolate
causal effects of the teaching mechanism.}

DPG achieves fewer learner updates within the time budget. Selfplay moves much
further from the base sampling distribution. The fixed hard reference also has
nonzero KL because its distribution differs from the base by construction.

\clearpage
\section*{6. Interpretation limits and reproducibility}
\textbf{Single-seed scope.} All scientific runs use seed 20261017; there is no
independent-seed uncertainty estimate. The tuning scope is narrow (generator
learning rate 0.001). These synthetic experiments use a finite, 12-regime Heston
menu. They do not establish performance on other distributions or live markets.

\textbf{Small hard-domain differences.} The archived report describes a separate
32-episode-per-domain pricing sensitivity check. No baseline-relative signs
reversed in that bounded check, but it does not certify the full-test gaps or
rare-tail behavior. Small hard-domain improvements should therefore be treated
as measured differences rather than robust superiority.

\textbf{Loss is not the whole financial outcome.} MSE does not summarize all
signed PnL, expected shortfall, turnover or transaction-cost behavior. This
report concentrates on the requested training and evaluation losses; the source
CSVs retain additional metrics.

\textbf{Time comparisons.} Runs use the same GPU model but not controlled hosts
or clocks. Active time excludes evaluation, compilation and I/O. Equal active
training time is neither equal learner-step count nor equal total compute cost.

\subsection*{Local reproducibility}
The report directory includes the LaTeX source, vector PDF figures, PNG previews,
the plotting script, archived CSV inputs and a SHA-256 manifest. Plot inputs are
\texttt{all\_events.csv}; the selected-results table comes from
\texttt{final\_metrics.csv}. Frozen identifiers are read from
\texttt{frozen\_selection.json}. The archived \texttt{report.md} supplies campaign
context. A snapshot of the relevant runner, objective and metric source files
records the loss definitions inspected for this report.

The build script checks finite positive losses, unique evaluation records and
agreement between each frozen selected checkpoint's final metrics and its raw
event rows. It filters scientific runs using the frozen-selection manifest,
sorts curves by cumulative active time, and removes duplicate training records
at the same run/update count. It does not recompute GPU metrics from episode
arrays or rerun training.

From the report directory, rebuild with:
\begin{verbatim}
python build_report.py
latexmk -pdf -interaction=nonstopmode -halt-on-error report.tex
\end{verbatim}
Python dependencies: NumPy, pandas and Matplotlib. A standard LaTeX installation
with the packages declared in \texttt{report.tex} is required. The archived
\texttt{manifest.json} identifies the exact local source snapshot used.

\end{document}
